\chapter{C4 Architecture Modeling}
\label{ch:c4-architecture}

\begin{chaptergoalsbox}
\textbf{What You Will Learn:}
\begin{itemize}[leftmargin=12pt]
  \item Master the four hierarchical abstraction levels of the C4 model: System Context (L1), Container (L2), Component (L3), and Code (L4).
  \item Apply strict visual syntax rules to eliminate diagram ambiguity: explicit element labeling [Name, Role, Technology], unidirectional labeled interaction arrows, and visual boundary scoping.
  \item Avoid common architectural communication failures, including the dreaded ``boxes and lines without semantics'' anti-pattern and the conflation of deployable software units with operating system virtualization.
  \item Implement Architecture-as-Code (AaC) workflows using the Structurizr Domain-Specific Language (DSL) to version, validate, and automatically render architectural models alongside source code.
\end{itemize}

\textbf{What You Should Already Understand:}
\begin{itemize}[leftmargin=12pt]
  \item Foundational distributed systems concepts: microservices, client-server models, message brokers, and relational or document datastores.
  \item Inter-service communication protocols including HTTP/REST, gRPC over HTTP/2, and asynchronous event streaming (e.g., Apache Kafka, RabbitMQ).
  \item Standard software development lifecycles, Git revision control, and Continuous Integration (CI) delivery pipelines.
\end{itemize}

\textbf{Why This Topic Exists:}
Software systems grow rapidly in complexity. Traditional Unified Modeling Language (UML) specifications often present steep learning curves and overwhelming syntactic ceremony, whereas ad-hoc whiteboard sketches lack rigor, consistency, and standard notation. The C4 model, developed by Simon Brown, provides an intuitive, hierarchical map of software systems that enables engineers, architects, and non-technical stakeholders to communicate architecture at varying levels of zoom without losing semantic clarity.

\textbf{Motivating Practical Problem:}
A high-growth fintech platform expands from a single monolithic application into 45 microservices across multiple cloud regions. During a Sev-1 production outage involving settlement failures, the on-call incident team discovers that no single engineer understands the end-to-end payment flow. Existing architecture diagrams in the company wiki consist of 3-year-old, static PNG exports with unlabelled arrows, obsolete services, and conflicting boundaries. The team loses four hours simply identifying which service publishes payment events to the clearinghouse.
\end{chaptergoalsbox}

\section{The Communication Crisis in Software Architecture}

Software architecture is fundamentally about communication, decision-making, and managing complexity. Yet, across the global software industry, architectural documentation routinely fails in practice. When engineers are asked to draw their system on a whiteboard, they almost invariably produce a confusing maze of rectangles, unlabeled circles, and ambiguous bidirectional arrows. 

Consider the typical diagramming anti-patterns encountered in enterprise engineering:
\begin{enumerate}[leftmargin=18pt]
  \item \textbf{The Mysterious Arrow}: A line connecting ``Service A'' to ``Service B'' with an arrowhead on both ends and no text label. Does this represent a synchronous gRPC call, a shared database table, a WebSocket duplex stream, or data flow direction?
  \item \textbf{The Missing Legend}: Color-coded nodes (blue boxes, red boxes, green cylinders) without an explanatory key, forcing readers to guess whether colors signify programming languages, cloud providers, or security domains.
  \item \textbf{Mixed Levels of Abstraction}: A single diagram that depicts a high-level customer persona, an AWS Availability Zone, an individual React component, and a PostgreSQL table side by side.
  \item \textbf{Acronym Overload and Ghost Services}: Services identified only by internal project codenames (e.g., ``Hydra'' or ``Apollo'') with zero indication of their functional responsibility or technical implementation.
\end{enumerate}

To solve this crisis, software architect Simon Brown created the **C4 model**. The core insight of C4 is that software architecture documentation should function like a geographical map service, such as Google Maps. When viewing the entire planet, you see continental boundaries and major nations (System Context). As you zoom into a metropolitan region, major highways and transit networks appear (Containers). Zooming further exposes individual streets and intersections (Components). Finally, street-level imagery reveals building doorways and house numbers (Code).

\begin{intuitionbox}[The Geographical Map Analogy]
No sensible cartographer attempts to display every residential driveway on a global wall map. Similarly, effective software architects never show database tables or class inheritance trees on an enterprise landscape diagram. The C4 model enforces strict abstraction boundaries: you select the zoom level appropriate for your audience and technical purpose, keeping cognitive load within human working memory limits.
\end{intuitionbox}

\section{The Four Hierarchical Abstraction Levels}

The C4 model structures software architecture into four distinct, nested tiers: \textbf{Context}, \textbf{Containers}, \textbf{Components}, and \textbf{Code}. Each tier expands upon an element defined in the layer above it.

\begin{figure}[htbp]
\centering
\begin{tikzpicture}[
  node distance=1.0cm and 1.2cm,
  levelbox/.style={rectangle, draw=headingblue!80, fill=headingblue!5, rounded corners=3pt, text width=12cm, minimum height=1.3cm, inner sep=8pt, font=\sffamily},
  titlebox/.style={font=\sffamily\bfseries\color{headingblue}},
  descbox/.style={font=\sffamily\small\color{textdark}}
]

\node[levelbox] (l1) {
  \textbf{\color{headingblue}Level 1: System Context Diagram}\\
  \small\color{textdark}Defines the software system within its environment. Shows human users (personas) and external software systems with which it interacts. High-level scope boundary; technology-agnostic.
};

\node[levelbox, below=0.6cm of l1] (l2) {
  \textbf{\color{headingblue}Level 2: Container Diagram}\\
  \small\color{textdark}Decomposes the system into deployable/executable units: web applications, mobile apps, API services, serverless workers, databases, file stores. Identifies core technologies and inter-container protocols.
};

\node[levelbox, below=0.6cm of l2] (l3) {
  \textbf{\color{headingblue}Level 3: Component Diagram}\\
  \small\color{textdark}Zooms into an individual container to reveal its internal modular architecture: controllers, service facades, repositories, domain modules, and internal dependencies.
};

\node[levelbox, below=0.6cm of l3] (l4) {
  \textbf{\color{headingblue}Level 4: Code Diagram (Optional)}\\
  \small\color{textdark}Shows implementation details: UML class diagrams, database schemas, or entity-relationship models. Typically generated automatically via IDEs; rarely maintained manually.
};

\draw[-{Stealth[scale=1.1]}, thick, color=headingblue!60] (l1) -- (l2);
\draw[-{Stealth[scale=1.1]}, thick, color=headingblue!60] (l2) -- (l3);
\draw[-{Stealth[scale=1.1]}, thick, color=headingblue!60] (l3) -- (l4);

\end{tikzpicture}
\caption{The C4 Model Hierarchical Zoom Progression}
\label{fig:c4-hierarchy}
\end{figure}

\subsection{Level 1: System Context Diagram}

The System Context diagram sets the stage. It establishes the boundary between what we are building (the system under discussion) and the outside world. It answers three fundamental questions:
\begin{itemize}[leftmargin=14pt]
  \item Who are the primary human actors (customers, support staff, administrators) utilizing the system?
  \item What high-level software system is being delivered?
  \item What external software systems, third-party APIs, or legacy mainframes does this system integrate with?
\end{itemize}

At this level, internal implementation details, such as programming languages, database engines, and cloud hosting platforms, are intentionally omitted. The audience includes everyone from product managers and C-suite executives to junior engineers.

\begin{figure}[htbp]
\centering
\begin{tikzpicture}[
  node distance=1.5cm and 2.0cm,
  person/.style={rectangle, draw=gray!70, fill=gray!15, rounded corners=4pt, text width=3.0cm, minimum height=1.6cm, align=center, inner sep=6pt, font=\sffamily\footnotesize},
  system/.style={rectangle, draw=exampleBlue!80, fill=exampleBg, rounded corners=4pt, text width=3.8cm, minimum height=1.8cm, align=center, inner sep=6pt, font=\sffamily\footnotesize},
  extsystem/.style={rectangle, draw=gray!60, fill=gray!10, rounded corners=4pt, text width=3.4cm, minimum height=1.8cm, align=center, inner sep=6pt, font=\sffamily\footnotesize},
  boundary/.style={draw=exampleBlue!50, dashed, rounded corners=6pt, inner sep=12pt},
  arrow/.style={-{Stealth[scale=1.0]}, thick, color=textdark}
]

\node[person] (merchant) {\textbf{Merchant}\\ \scriptsize[Person]\\ \raggedright Sells goods online, manages products, reviews settlements.};
\node[system, right=2.0cm of merchant] (paycore) {\textbf{Global Payments Platform}\\ \scriptsize[Software System]\\ \raggedright Accepts card transactions, settles multi-currency funds, manages escrow.};
\node[extsystem, right=2.0cm of paycore] (bank) {\textbf{Banking Rails / ACH}\\ \scriptsize[External System]\\ \raggedright Automated Clearing House for inter-bank bulk transfers.};
\node[extsystem, above=1.5cm of paycore] (fraud) {\textbf{Third-Party Fraud API}\\ \scriptsize[External System]\\ \raggedright Real-time AML and credit risk verification provider.};
\node[person, below=1.5cm of merchant] (buyer) {\textbf{Shopper}\\ \scriptsize[Person]\\ \raggedright Purchases products via checkout widgets.};

\draw[arrow] (merchant) -- node[above, font=\tiny\sffamily, align=center]{Manages account\\via HTTPS} (paycore);
\draw[arrow] (buyer) |- node[below, pos=0.7, font=\tiny\sffamily, align=center]{Submits payments\\via HTTPS} (paycore);
\draw[arrow] (paycore) -- node[above, font=\tiny\sffamily, align=center]{Transfers funds\\via ISO 20022} (bank);
\draw[arrow] (paycore) -- node[right, font=\tiny\sffamily, align=center]{Scores transactions\\via gRPC/mTLS} (fraud);

\end{tikzpicture}
\caption{C4 Level 1: System Context Diagram for Global Payments Platform}
\label{fig:c4-l1-context}
\end{figure}

\subsection{Level 2: Container Diagram}

In C4 terminology, a \textbf{container} is NOT a Docker container or Kubernetes pod. A container is defined as any \textit{separately deployable or runnable unit of software that executes code or stores data}. 

Typical examples of C4 containers include:
\begin{itemize}[leftmargin=14pt]
  \item Single-Page Applications (e.g., React, Vue, Svelte bundle served from an S3 bucket or CDN).
  \item Mobile Applications (e.g., iOS Swift, Android Kotlin, React Native binary running on a client device).
  \item Backend Web Services and APIs (e.g., Node.js Express service, Go HTTP microservice, Python FastAPI worker).
  \item Serverless Functions (e.g., AWS Lambda, Google Cloud Run service).
  \item Persistent Data Stores (e.g., PostgreSQL relational database, MongoDB document store, DynamoDB).
  \item In-Memory Caches and Message Brokers (e.g., Redis cluster, Apache Kafka broker, RabbitMQ exchange).
  \item File Storage Systems (e.g., AWS S3, Azure Blob Storage).
\end{itemize}

The Container diagram zooms inside the software system boundary established in Level 1. It details the high-level technical choices, responsibilities, and network communication protocols (e.g., HTTPS, gRPC, WebSocket, TCP, JDBC).

\begin{figure}[htbp]
\centering
\begin{tikzpicture}[
  node distance=1.2cm and 1.4cm,
  cnt/.style={rectangle, draw=headingblue!70, fill=headingblue!8, rounded corners=3pt, text width=3.4cm, minimum height=1.6cm, align=left, inner sep=6pt, font=\sffamily\scriptsize},
  db/.style={cylinder, shape border rotate=90, aspect=0.25, draw=headingblue!70, fill=headingblue!8, text width=2.8cm, minimum height=1.5cm, align=center, inner sep=4pt, font=\sffamily\scriptsize},
  arrow/.style={-{Stealth[scale=0.9]}, thick, color=textdark}
]

% Boundary
\draw[draw=gray!50, dashed, rounded corners=6pt] (-0.3,-5.2) rectangle (12.8,1.8);
\node[font=\sffamily\bfseries\color{gray!80}, anchor=north west] at (0.0,1.6) {Global Payments Platform [System Boundary]};

% Nodes
\node[cnt] (web) at (1.5, 0.4) {\textbf{Merchant Web App}\\ \textit{[Container: React / TypeScript]}\\ Delivers dashboard UI in browser.};
\node[cnt] (api) at (6.5, 0.4) {\textbf{API Gateway}\\ \textit{[Container: Go / Envoy]}\\ Routing, auth verification, and rate limiting.};
\node[cnt] (order) at (6.5, -2.2) {\textbf{Settlement Service}\\ \textit{[Container: Java / Spring Boot]}\\ Orchestrates ledger balances and payout scheduling.};
\node[db] (sqldb) at (11.0, -2.2) {\textbf{Ledger DB}\\ \textit{[Container: PostgreSQL]}\\ Multi-currency double-entry ledger.};
\node[cnt] (kafka) at (1.5, -4.0) {\textbf{Event Broker}\\ \textit{[Container: Apache Kafka]}\\ Distributed event streaming log for settlements.};
\node[cnt] (worker) at (6.5, -4.0) {\textbf{Payout Worker}\\ \textit{[Container: Python AsyncIO]}\\ Pulls batch payout events, interacts with ACH.};

% Connections
\draw[arrow] (web) -- node[above, font=\tiny\sffamily]{Calls API [JSON/HTTPS]} (api);
\draw[arrow] (api) -- node[right, font=\tiny\sffamily]{Forwards [gRPC]} (order);
\draw[arrow] (order) -- node[above, font=\tiny\sffamily]{Read/Write [JDBC/TLS]} (sqldb);
\draw[arrow] (order) -- node[left, font=\tiny\sffamily]{Publishes [Kafka Wire]} (kafka);
\draw[arrow] (kafka) -- node[above, font=\tiny\sffamily]{Consumes} (worker);

\end{tikzpicture}
\caption{C4 Level 2: Container Diagram for Settlement Architecture}
\label{fig:c4-l2-container}
\end{figure}

\begin{commonmistakebox}[Equating C4 Containers to OCI / Docker Containers]
The word ``container'' in the C4 model was coined before Docker became the ubiquitous industry standard for packaging software. A C4 container is an \textit{architectural deployment boundary}. A PostgreSQL database installed natively on bare-metal Linux is a C4 container. A client-side iOS Swift application running on an iPhone is a C4 container. Do not confuse the C4 abstraction with an OCI container image or a Kubernetes pod.
\end{commonmistakebox}

\subsection{Level 3: Component Diagram}

The Component diagram zooms inside an individual container from Level 2 to show its internal logical components. Simon Brown defines a component as: \textit{``A grouping of related functionality encapsulated behind a well-defined interface, running inside a container.''}

In modern enterprise architectures, components reflect modular subsystems:
\begin{itemize}[leftmargin=14pt]
  \item Controllers or API Facades handling HTTP request deserialization and schema validation.
  \item Domain Services implementing core business workflows and transaction boundaries.
  \item Repositories or Data Access Objects (DAOs) abstracting database queries and persistence logic.
  \item External Client Adapters wrapping third-party SDKs, circuit breakers, and rate limiters.
\end{itemize}

The primary consumers of Component diagrams are software engineers actively working on that specific codebase. Component diagrams illustrate separation of concerns, dependency injection graphs, and architectural compliance rules (e.g., verifying that a controller never calls a database driver directly, bypassing domain invariants).

\subsection{Level 4: Code Diagram}

Level 4 zooms into an individual component to show how it is implemented as code artifacts. This typically takes the form of UML Class diagrams or Entity-Relationship diagrams. 

In practice, manual creation and maintenance of Level 4 diagrams is an anti-pattern. Code changes daily; maintaining handwritten class diagrams produces immediate documentation staleness. Level 4 should only be generated on-demand by automated tooling (such as IDE diagram generators, Doxygen, or TypeORM schema exporters) for uniquely complex design patterns or algorithmic state machines.

\section{Visual Syntax Rules and Anti-Patterns}

To ensure diagrams communicate unambiguously, the C4 model establishes explicit visual notation conventions:

\begin{table}[htbp]
\centering\small
\begin{tabular}{lp{5.5cm}p{5.5cm}}
\toprule
\textbf{Element} & \textbf{Mandatory Requirement} & \textbf{Anti-Pattern to Avoid} \\
\midrule
\textbf{Box Title} & Explicit name + element role (e.g., \texttt{Ledger Service [Container: Go]}). & Cryptic internal codename only (e.g., \texttt{Zeus}). \\
\textbf{Technology Tag} & Specific implementation language, framework, or vendor. & Blank technology box leaving readers to guess. \\
\textbf{Description} & One or two sentences summarizing primary responsibility. & Empty boxes with only a title. \\
\textbf{Arrows} & Unidirectional; describes the intent of the trigger. & Bidirectional arrows without flow semantics. \\
\textbf{Arrow Labels} & States action and wire protocol (e.g., \texttt{Publishes balance update to [Kafka]}). & Unlabeled line or generic label like ``Uses'' or ``Data''. \\
\textbf{Boundaries} & Dashed box with system or enterprise boundary title. & Floating elements without clear architectural ownership. \\
\bottomrule
\end{tabular}
\caption{C4 Visual Syntax Checklist and Anti-Patterns}
\label{tab:c4-syntax-checklist}
\end{table}

\begin{productionwarningbox}[Architecture Drift and Static Wiki Rot]
Static PNG and Visio diagrams stored in corporate wikis degrade rapidly over time. Within three months of a team restructuring or major release, static diagrams become misleading liability artifacts. In regulated industries (e.g., SOC2, PCI-DSS, ISO 27001), presenting outdated architecture diagrams to auditors during an incident investigation can result in immediate compliance failure. Architecture must be treated with the same engineering rigor as code.
\end{productionwarningbox}

\section{Architecture-as-Code: The Structurizr DSL}

To maintain architectural documentation alongside source code, modern engineering teams utilize \textbf{Architecture-as-Code (AaC)}. The gold standard for C4 modeling is the **Structurizr DSL**, designed by Simon Brown. 

The Structurizr DSL decouples the \textit{architectural model} (the graph of people, software systems, containers, components, and relationships) from the \textit{views} (visual renderings, layouts, and styles).

\begin{lstlisting}[language=Java, caption={Complete Structurizr DSL Definition for an E-Commerce Payment Gateway}]
workspace "Global Payments Gateway" "Production C4 Model" {

    model {
        merchant = person "Online Merchant" "An enterprise merchant managing checkouts."
        shopper = person "End Shopper" "A customer purchasing goods online."

        enterprise "Acme Fintech Corp" {
            fraudApi = softwareSystem "Risk Engine" "Scores transactions for credit fraud." "External"
            bankingRails = softwareSystem "ACH Settlement Rails" "Federal Reserve clearinghouse." "External"

            paymentSystem = softwareSystem "Payment Gateway" "Processes global card payments." {
                webApp = container "Merchant Portal" "Next.js dashboard for settlements" "TypeScript / React"
                apiGateway = container "Edge Gateway" "TLS termination, auth, and rate limiting" "Envoy / Go"
                paymentCore = container "Transaction Engine" "Executes state machine for payments" "Java / Spring Boot"
                ledgerDb = container "Ledger Store" "Immutable double-entry transaction store" "PostgreSQL 16" "Database"
                eventBus = container "Payment Bus" "High-throughput event streaming" "Apache Kafka" "Queue"
            }
        }

        # User interactions
        shopper -> apiGateway "Submits checkout payments" "JSON/HTTPS"
        merchant -> webApp "Views settlement reports" "HTTPS"
        webApp -> apiGateway "Calls backend endpoints" "JSON/HTTPS"

        # Internal container interactions
        apiGateway -> paymentCore "Dispatches charge requests" "gRPC/HTTP2"
        paymentCore -> ledgerDb "Appends journal entries" "JDBC/mTLS"
        paymentCore -> eventBus "Emits PaymentCaptured events" "Kafka Wire Protocol"
        paymentCore -> fraudApi "Requests fraud risk score" "JSON/HTTPS"
        paymentCore -> bankingRails "Submits batch settlement file" "SFTP / ISO 20022"
    }

    views {
        systemContext paymentSystem "SystemContext" {
            include *
            autoLayout lr
        }

        container paymentSystem "Containers" {
            include *
            autoLayout tb
        }

        styles {
            element "Person" {
                shape Person
                background #08427b
                color #ffffff
            }
            element "Software System" {
                background #1168bd
                color #ffffff
            }
            element "Container" {
                background #438dd5
                color #ffffff
            }
            element "Database" {
                shape Cylinder
                background #2a52be
                color #ffffff
            }
        }
    }
}
\end{lstlisting}

\begin{underthehoodbox}[How Structurizr Compiles Architecture Graphs]
When the Structurizr CLI processes a `.dsl` file, it constructs an in-memory Abstract Syntax Tree (AST) of the workspace. It validates referential integrity (e.g., verifying that relationships only reference declared elements) and exports a canonical JSON representation conforming to the Structurizr Schema. From this JSON, headless renderers can automatically generate PlantUML, Mermaid.js, WebGL diagrams, or interactive SVG views inside documentation sites, ensuring diagrams are version-controlled in Git and re-rendered on every CI commit.
\end{underthehoodbox}

\begin{realworldbox}[Managing 200+ Microservices at Monzo Bank]
Digital bank Monzo famously manages hundreds of microservices authored in Go. Rather than drawing manual diagrams, their platform team continuously harvests service metadata, RPC routing tables, and Kafka topic subscriptions directly from Git repos and Kubernetes manifests. This metadata is programmatically transformed into hierarchical C4 models. Any engineer can instantly generate an up-to-the-minute C4 container diagram centered on any microservice, showing its exact upstream callers, downstream dependencies, and storage backing.
\end{realworldbox}

\begin{tryityourselfbox}[Exercise: Model a Notification Service]
Using the Structurizr DSL syntax detailed above, model a multi-tenant Notification Service container that:
\begin{enumerate}
  \item Consumes notification requests from a shared Apache Kafka cluster.
  \item Dispatches push notifications to Apple APNs and Google FCM.
  \item Sends transactional emails via SendGrid.
  \item Tracks delivery status in a Redis cluster.
\end{enumerate}
Declare the containers, assign appropriate technology labels, and establish unidirectional relationships with wire protocols.
\end{tryityourselfbox}

\section{Summary \protect\& Key Takeaways}

\begin{itemize}[leftmargin=18pt]
  \item \textbf{Hierarchical Navigation}: The C4 model prevents abstraction confusion by partitioning architecture into four tiers: System Context (L1), Containers (L2), Components (L3), and Code (L4).
  \item \textbf{Precise Terminology}: In C4, a container is a deployable runtime unit (e.g., web server, database, mobile app), not specifically a Docker container.
  \item \textbf{Unambiguous Labeling}: Every element must specify its Name, Technology, and Responsibility. Every arrow must specify an action and communication protocol.
  \item \textbf{Avoid Code Diagram Rot}: Avoid manually maintaining Level 4 (Code) diagrams; leverage automated reflection or IDE tools when class-level views are required.
  \item \textbf{Architecture-as-Code}: Maintain models in text-based DSLs (such as Structurizr DSL) inside Git repositories to enable peer review, versioning, and automated rendering.
\end{itemize}

\section{Exercises}

\subsection*{Mechanics}
\begin{enumerate}[leftmargin=18pt]
  \item Given the following system components: (a) PostgreSQL 15, (b) React Web Client, (c) Customer Support Specialist, (d) Stripe Billing API, (e) Django REST Framework backend. Classify each element into its proper C4 abstraction category (Person, Software System, External System, or Container).
  \item Explain why a bidirectional arrow without labels between an API service and a Redis cache violates C4 diagramming standards. Re-draw or describe the relationship using proper C4 conventions.
\end{enumerate}

\subsection*{Debugging}
\begin{enumerate}[leftmargin=18pt]
  \item Review the following snippet from an engineering wiki: \textit{``Our System Context diagram shows the User communicating directly with our Python SQLAlchemy ORM Session, which bypasses the Nginx Reverse Proxy.''} Identify two severe abstraction-level violations in this statement.
\end{enumerate}

\subsection*{Architecture Design}
\begin{enumerate}[leftmargin=18pt]
  \item Design a C4 Level 2 Container Diagram for a high-throughput IoT Telemetry Ingestion Platform. The platform receives MQTT packets from 100,000 smart energy meters, ingests them via an EMQX broker, processes readings using an Apache Flink stream processor, stores time-series data in TimescaleDB, and exposes an internal Go REST API for analytical dashboards.
\end{enumerate}

\section{Detailed Solutions}

\subsection*{Mechanics Solution}
\begin{enumerate}[leftmargin=18pt]
  \item 
  \begin{itemize}
    \item (a) \textbf{PostgreSQL 15}: Container (Data Store).
    \item (b) \textbf{React Web Client}: Container (Single-Page Application).
    \item (c) \textbf{Customer Support Specialist}: Person.
    \item (d) \textbf{Stripe Billing API}: External Software System.
    \item (e) \textbf{Django REST Framework backend}: Container (API Application).
  \end{itemize}
  \item A bidirectional arrow fails to communicate causality, request direction, and network protocol. In reality, the API Service initiates a TCP connection to Redis to query or store keys. The corrected relationship is a unidirectional arrow:
  $$\text{API Service} \xrightarrow{\text{Reads and writes session tokens via Redis Wire Protocol [TCP/6379]}} \text{Redis Cluster}$$
\end{enumerate}

\subsection*{Debugging Solution}
\begin{enumerate}[leftmargin=18pt]
  \item The statement commits two major abstraction violations:
  \begin{enumerate}
    \item \textbf{Level Mixing (L1 vs L3/L4)}: A System Context diagram (L1) must only depict high-level systems and persons. An ORM session (SQLAlchemy) is a Level 3 Component or Level 4 Code concept; it has no place in a System Context view.
    \item \textbf{Violation of Network Topologies}: A human user in a browser cannot open an in-memory Python SQLAlchemy ORM session directly over the public internet. The user communicates with an edge proxy or web container over HTTPS.
  \end{enumerate}
\end{enumerate}

\subsection*{Architecture Design Solution}
\begin{figure}[htbp]
\centering
\begin{tikzpicture}[
  node distance=1.2cm and 1.5cm,
  cnt/.style={rectangle, draw=headingblue!80, fill=headingblue!8, rounded corners=3pt, text width=3.3cm, minimum height=1.6cm, align=left, inner sep=6pt, font=\sffamily\scriptsize},
  db/.style={cylinder, shape border rotate=90, aspect=0.25, draw=headingblue!80, fill=headingblue!8, text width=2.8cm, minimum height=1.5cm, align=center, inner sep=4pt, font=\sffamily\scriptsize},
  arrow/.style={-{Stealth[scale=0.9]}, thick, color=textdark}
]

\node[cnt] (meter) at (0, 0) {\textbf{Smart Energy Meter}\\ \textit{[External System]}\\ Transmits telemetry telemetry.};
\node[cnt] (broker) at (4.5, 0) {\textbf{MQTT Broker}\\ \textit{[Container: EMQX]}\\ Ingests telemetry streams.};
\node[cnt] (flink) at (9.0, 0) {\textbf{Stream Pipeline}\\ \textit{[Container: Apache Flink]}\\ Aggregates 1-minute window metrics.};
\node[db] (tsdb) at (9.0, -2.8) {\textbf{Telemetry Store}\\ \textit{[Container: TimescaleDB]}\\ Hypertable partition storage.};
\node[cnt] (api) at (4.5, -2.8) {\textbf{Analytics API}\\ \textit{[Container: Go / Gin]}\\ Serves aggregation metrics.};
\node[cnt] (ui) at (0, -2.8) {\textbf{Grid Ops Portal}\\ \textit{[Container: React / TypeScript]}\\ Real-time grid dashboard.};

\draw[arrow] (meter) -- node[above, font=\tiny\sffamily]{Sends telemetry [MQTT/TLS]} (broker);
\draw[arrow] (broker) -- node[above, font=\tiny\sffamily]{Consumes [Kafka protocol]} (flink);
\draw[arrow] (flink) -- node[right, font=\tiny\sffamily]{Writes [JDBC]} (tsdb);
\draw[arrow] (api) -- node[above, font=\tiny\sffamily]{Queries [SQL/TCP]} (tsdb);
\draw[arrow] (ui) -- node[above, font=\tiny\sffamily]{Fetches metrics [HTTPS/JSON]} (api);

\end{tikzpicture}
\caption{C4 Level 2 Container Model for IoT Ingestion Platform}
\end{figure}

\section{Further Reading \protect\& Citations}
\begin{itemize}[leftmargin=18pt]
  \item Brown, S. (2018). \textit{Software Architecture for Developers: Visualise, document and explore your software architecture}. Leanpub.
  \item Kruchten, P. (1995). Architectural Blueprints---The ``4+1'' View Model of Software Architecture. \textit{IEEE Software}, 12(6), 42--50.
  \item Structurizr DSL Specification. Available at: \url{https://structurizr.com/dsl}.
  \item IEEE Computer Society. (2000). \textit{IEEE Recommended Practice for Architectural Description of Software-Intensive Systems} (IEEE Std 1471-2000).
\end{itemize}
